\chapter{Aim and Scope of the Thesis}

\section{Main Aim}

The main aim of this thesis is to develop a nonlinear simulation environment for a Boeing 777-like transport aircraft and to design an autopilot concept based on Nonlinear Model Predictive Control. The simulator is intended to support analysis of six-degree-of-freedom aircraft motion, trimming in representative flight conditions and evaluation of constrained control algorithms.

\section{Detailed Objectives}

The detailed objectives of the work are:
\begin{enumerate}
    \item to define coordinate frames, sign conventions, state vectors and input vectors used throughout the project,
    \item to prepare a first-cut geometry, mass and inertia model of a Boeing 777-like aircraft,
    \item to implement a simple International Standard Atmosphere model and fixed project constants,
    \item to implement a simplified aerodynamic coefficient model suitable for early flight-dynamics integration,
    \item to define a propulsion and actuator modelling approach,
    \item to formulate the nonlinear six-degree-of-freedom equations of motion,
    \item to prepare trimming and validation procedures for selected reference cases,
    \item to design longitudinal and lateral-directional NMPC control structures,
    \item to evaluate the model and controller using repeatable simulation scenarios.
\end{enumerate}

\section{Scope}

The scope of the thesis includes rigid-body flight dynamics, simplified aerodynamics, simplified propulsion, actuator limits, standard atmosphere modelling, trimming, qualitative validation and NMPC autopilot design. The implementation target is MATLAB/Simulink, with MATLAB data modules used to define reusable aircraft and environment parameters.

The following topics are outside the scope of the present work:
\begin{itemize}
    \item structural flexibility and aeroelastic effects,
    \item detailed hydraulic, electrical or avionics system models,
    \item certified flight-control laws,
    \item proprietary aerodynamic or engine data,
    \item high-fidelity pilot-display or flight-management-system simulation,
    \item certification-level validation.
\end{itemize}

\section{Assumptions and Limitations}

The aircraft is treated as a rigid body. The Earth is initially approximated as flat and non-rotating, and the navigation frame is a local North-East-Down frame. All internal calculations are performed in SI units. The current implementation uses fixed standard gravity $g_0=\SI{9.80665}{\metre\per\square\second}$ as a project-wide constant. Gravity-dependent force quantities are computed from mass and the selected gravity value rather than stored as immutable mass properties.

The aerodynamic model is intentionally simple. Its purpose is to provide physically reasonable signs, magnitudes and coupling terms for early simulation and control-design studies. The model is expected to be refined after trim tests and dynamic-response validation.

\section{Success Criteria}

The project will be considered successful if the following criteria are met:
\begin{enumerate}
    \item the six-degree-of-freedom plant can be executed in simulation without numerical instability caused by modelling errors,
    \item trim points can be obtained for several representative flight conditions,
    \item open-loop dynamic responses are qualitatively consistent with expectations for a large transport aircraft,
    \item the NMPC controller can stabilise and track selected flight variables,
    \item actuator, thrust and selected state constraints are respected during simulation,
    \item the obtained results are documented with plots, tables and engineering interpretation.
\end{enumerate}
